\documentclass[12pt,a4paper]{article}
\usepackage[utf8]{inputenc}
\usepackage[vietnamese]{babel}
\usepackage{amsmath,amssymb,amsthm}
\usepackage{geometry}
\usepackage{enumitem}
\usepackage[most]{tcolorbox}
\usepackage{hyperref}
\usepackage{fancyhdr}
\usepackage{tikz}
\usepackage{eso-pic}
\usepackage{xcolor}
\geometry{a4paper, margin=2cm, bottom=2.5cm}
\hypersetup{colorlinks=true, linkcolor=blue!70!black}
\setlist[itemize]{topsep=4pt, itemsep=2pt}
\setlist[enumerate]{topsep=4pt, itemsep=2pt}
\AddToShipoutPictureFG{%
  \AtPageCenter{%
    \begin{tikzpicture}[remember picture, overlay]
      \node[rotate=45, scale=1, opacity=0.06]{\fontsize{60}{72}\selectfont\bfseries\sffamily Đặng Tấn Quí};
    \end{tikzpicture}%
  }%
}
\pagestyle{fancy}
\fancyhf{}
\fancyhead[L]{\small\textit{GIẢI TÍCH 3}}
\fancyhead[R]{\small\textit{Đặng Tấn Quí}}
\fancyfoot[C]{\thepage}
\fancyfoot[L]{\footnotesize\textcopyright\ Đặng Tấn Quí}
\fancyfoot[R]{\footnotesize SĐT: 0377\,122\,210}
\renewcommand{\headrulewidth}{0.4pt}
\renewcommand{\footrulewidth}{0.4pt}
\fancypagestyle{plain}{\fancyhf{}\fancyfoot[C]{\thepage}\fancyfoot[L]{\footnotesize\textcopyright\ Đặng Tấn Quí}\fancyfoot[R]{\footnotesize SĐT: 0377\,122\,210}\renewcommand{\headrulewidth}{0pt}\renewcommand{\footrulewidth}{0.4pt}}
\newtcolorbox{dinhnghia}[1][]{enhanced, breakable, colback=blue!3!white, colframe=blue!60!black, fonttitle=\bfseries, title={Định nghĩa}, #1}
\newtcolorbox{dinhly}[1][]{enhanced, breakable, colback=green!3!white, colframe=green!50!black, fonttitle=\bfseries, title={Định lý}, #1}
\newtcolorbox{tinhchat}[1][]{enhanced, breakable, colback=yellow!5!white, colframe=orange!50!black, fonttitle=\bfseries, title={Tính chất}, #1}
\newtcolorbox{phuongphap}[1][]{enhanced, breakable, colback=gray!5!white, colframe=gray!50!black, fonttitle=\bfseries, title={Phương pháp}, #1}
\newtcolorbox{luuy}[1][]{enhanced, breakable, colback=red!3!white, colframe=red!50!black, fonttitle=\bfseries, title={Lưu ý}, #1}
\newtcolorbox{congthuc}[1][]{enhanced, breakable, colback=cyan!3!white, colframe=cyan!50!black, fonttitle=\bfseries, title={Công thức}, #1}
\newtcolorbox{vidu}[1][]{enhanced, breakable, colback=violet!3!white, colframe=violet!50!black, fonttitle=\bfseries, title={Ví dụ}, #1}

\begin{document}
\thispagestyle{empty}
\begin{tikzpicture}[remember picture, overlay]
  \fill[blue!80!black] (current page.south west) rectangle (current page.north east);
  \node[anchor=west, white, font=\fontsize{26}{32}\bfseries\selectfont]
    at ([xshift=3cm, yshift=-7.5cm]current page.north west) {GIẢI TÍCH 3};
  \node[anchor=west, white, font=\fontsize{18}{24}\selectfont]
    at ([xshift=3cm, yshift=-9cm]current page.north west) {Chuỗi số, chuỗi hàm, phương trình vi phân --- Năm 2};
  \node[anchor=west, white, font=\Large\bfseries] at ([xshift=3cm, yshift=-13cm]current page.north west) {Biên soạn:};
  \node[anchor=west, yellow!80!white, font=\fontsize{22}{28}\bfseries\selectfont] at ([xshift=3cm, yshift=-14.5cm]current page.north west) {ĐẶNG TẤN QUÍ};
  \node[anchor=south east, white, font=\Large\bfseries] at ([xshift=-2cm, yshift=3cm]current page.south east) {\the\year};
\end{tikzpicture}
\newpage
\tableofcontents
\newpage
%% ================================================================
%%  CHƯƠNG 1: CHUỖI SỐ
%% ================================================================
\section{Chuỗi số}

\begin{dinhnghia}
  \textbf{Chuỗi số} $\displaystyle\sum_{n=1}^{\infty} a_n$: tổng vô hạn $a_1 + a_2 + \cdots$. 
  
  Tổng riêng $S_n = a_1 + \cdots + a_n$. 
  
  Chuỗi \textbf{hội tụ} nếu $\displaystyle\lim_{n \to \infty} S_n$ tồn tại hữu hạn.
\end{dinhnghia}

\begin{dinhly}[title={Điều kiện cần hội tụ}]
  Nếu $\displaystyle\sum a_n$ hội tụ thì $\displaystyle\lim_{n \to \infty} a_n = 0$. (Điều ngược lại không đúng.)
\end{dinhly}

\begin{tinhchat}
  \begin{itemize}
    \item Nếu $\sum a_n$ hội tụ thì $\sum c\,a_n$ hội tụ ($c$ hằng số).
    \item Nếu $\sum a_n$ và $\sum b_n$ hội tụ thì $\sum (a_n \pm b_n)$ hội tụ.
    \item Bỏ hoặc thêm hữu hạn số hạng không ảnh hưởng tính hội tụ.
    \item Đặt ngoặc các số hạng liên tiếp không ảnh hưởng tổng (nếu chuỗi hội tụ).
  \end{itemize}
\end{tinhchat}

\subsection{Chuỗi số dương}

\begin{dinhly}[title={Tiêu chuẩn so sánh}]
  Cho $0 \le a_n \le b_n$ ($\forall n \ge n_0$):
  \begin{itemize}
    \item $\sum b_n$ hội tụ $\Rightarrow$ $\sum a_n$ hội tụ.
    \item $\sum a_n$ phân kỳ $\Rightarrow$ $\sum b_n$ phân kỳ.
  \end{itemize}
  \textbf{So sánh giới hạn:} Nếu $a_n, b_n > 0$ và $\displaystyle\lim \frac{a_n}{b_n} = c$ ($0 < c < +\infty$) thì $\sum a_n$ và $\sum b_n$ cùng hội tụ hoặc cùng phân kỳ.
\end{dinhly}

\begin{dinhly}[title={Tiêu chuẩn D'Alembert (tỉ số)}]
  Cho chuỗi dương $\sum a_n$, $\displaystyle\lim_{n \to \infty} \frac{a_{n+1}}{a_n} = L$:
  \begin{itemize}
    \item $L < 1$: chuỗi hội tụ.
    \item $L > 1$: chuỗi phân kỳ.
    \item $L = 1$: chưa kết luận.
  \end{itemize}
\end{dinhly}

\begin{dinhly}[title={Tiêu chuẩn Cauchy (căn)}]
  Cho chuỗi dương $\sum a_n$, $\displaystyle\lim_{n \to \infty} \sqrt[n]{a_n} = L$:
  \begin{itemize}
    \item $L < 1$: hội tụ. \quad $L > 1$: phân kỳ. \quad $L = 1$: chưa kết luận.
  \end{itemize}
\end{dinhly}

\begin{dinhly}[title={Tiêu chuẩn tích phân (Cauchy--Maclaurin)}]
  Nếu $f(x)$ liên tục, dương, giảm trên $[1, +\infty)$ và $a_n = f(n)$ thì $\sum_{n=1}^{\infty} a_n$ và $\int_1^{+\infty} f(x)\,dx$ cùng hội tụ hoặc cùng phân kỳ.
\end{dinhly}

\begin{congthuc}[title={Chuỗi mẫu}]
  \begin{itemize}
    \item \textbf{Chuỗi $p$:} $\displaystyle\sum_{n=1}^{\infty} \frac{1}{n^p}$: hội tụ khi $p > 1$, phân kỳ khi $p \le 1$.
    \item \textbf{Chuỗi hình học:} $\displaystyle\sum_{n=0}^{\infty} q^n = \frac{1}{1-q}$ khi $|q| < 1$; phân kỳ khi $|q| \ge 1$.
  \end{itemize}
\end{congthuc}

\subsection{Chuỗi có số hạng với dấu bất kỳ}

\begin{dinhnghia}
  \begin{itemize}
    \item $\sum a_n$ \textbf{hội tụ tuyệt đối} nếu $\sum |a_n|$ hội tụ.
    \item $\sum a_n$ \textbf{hội tụ bán phần} (hội tụ có điều kiện) nếu $\sum a_n$ hội tụ nhưng $\sum |a_n|$ phân kỳ.
  \end{itemize}
  Hội tụ tuyệt đối $\Rightarrow$ hội tụ. Điều ngược lại không đúng.
\end{dinhnghia}

\begin{dinhly}[title={Tiêu chuẩn Leibniz (chuỗi đan dấu)}]
  Chuỗi đan dấu $\displaystyle\sum_{n=1}^{\infty} (-1)^{n+1} a_n$ ($a_n > 0$) hội tụ nếu:
  \begin{enumerate}
    \item $(a_n)$ giảm: $a_{n+1} \le a_n$,
    \item $\displaystyle\lim_{n \to \infty} a_n = 0$.
  \end{enumerate}
  Khi đó tổng $S$ thỏa: $|S - S_n| \le a_{n+1}$.
\end{dinhly}

\begin{dinhly}[title={Tiêu chuẩn Abel và Dirichlet}]
  \begin{itemize}
    \item \textbf{Abel:} Nếu $\sum a_n$ hội tụ và $(b_n)$ đơn điệu bị chặn thì $\sum a_n b_n$ hội tụ.
    \item \textbf{Dirichlet:} Nếu dãy tổng riêng $\sum_{k=1}^n a_k$ bị chặn và $(b_n)$ đơn điệu giảm về $0$ thì $\sum a_n b_n$ hội tụ.
  \end{itemize}
\end{dinhly}

\subsection{Dãy hàm số}

\begin{dinhnghia}
  Dãy hàm $(f_n(x))$ trên $D$ \textbf{hội tụ điểm} đến $f(x)$ nếu $\forall x \in D$: $\displaystyle\lim_{n \to \infty} f_n(x) = f(x)$.

  Dãy hàm \textbf{hội tụ đều} đến $f$ trên $D$ nếu:
  \[
    \forall \varepsilon > 0,\; \exists N:\; \forall n > N,\; \forall x \in D:\; |f_n(x) - f(x)| < \varepsilon.
  \]
  Tương đương: $\displaystyle\lim_{n \to \infty} \sup_{x \in D} |f_n(x) - f(x)| = 0$.
\end{dinhnghia}

\begin{dinhly}[title={Tiêu chuẩn Cauchy về hội tụ đều}]
  $(f_n)$ hội tụ đều trên $D$ khi và chỉ khi $\forall \varepsilon > 0,\; \exists N:\; \forall m, n > N,\; \forall x \in D:\; |f_n(x) - f_m(x)| < \varepsilon$.
\end{dinhly}

\begin{tinhchat}[title={Tính chất dãy hàm hội tụ đều}]
  Nếu $(f_n)$ hội tụ đều đến $f$ trên $[a,b]$:
  \begin{enumerate}
    \item Nếu mỗi $f_n$ liên tục thì $f$ liên tục.
    \item $\displaystyle\lim_{n \to \infty} \int_a^b f_n(x)\,dx = \int_a^b f(x)\,dx$ (đổi giới hạn và tích phân).
    \item Nếu mỗi $f_n$ khả vi, $f_n'$ hội tụ đều thì $f' = \lim f_n'$ (đổi giới hạn và đạo hàm).
  \end{enumerate}
\end{tinhchat}

\subsection{Chuỗi hàm số}

\begin{dinhnghia}
  Chuỗi hàm $\displaystyle\sum_{n=1}^{\infty} u_n(x)$ \textbf{hội tụ} tại $x_0$ nếu chuỗi số $\sum u_n(x_0)$ hội tụ. Tập hợp các $x$ mà chuỗi hội tụ gọi là \textbf{miền hội tụ}.

  Chuỗi hàm \textbf{hội tụ đều} trên $D$ nếu dãy tổng riêng hội tụ đều trên $D$.
\end{dinhnghia}

\begin{dinhly}[title={Tiêu chuẩn Weierstrass (so sánh hội tụ đều)}]
  Nếu $|u_n(x)| \le M_n$ trên $D$ ($\forall n$) và $\sum M_n$ hội tụ thì $\sum u_n(x)$ hội tụ đều (và tuyệt đối) trên $D$.
\end{dinhly}

\begin{tinhchat}[title={Tính chất chuỗi hàm hội tụ đều}]
  Nếu $\sum u_n(x)$ hội tụ đều trên $[a,b]$:
  \begin{enumerate}
    \item Nếu mỗi $u_n$ liên tục thì tổng $S(x) = \sum u_n(x)$ liên tục.
    \item $\displaystyle\int_a^b S(x)\,dx = \sum_{n=1}^{\infty} \int_a^b u_n(x)\,dx$ (tích phân từng số hạng).
    \item Nếu mỗi $u_n$ khả vi, $\sum u_n'$ hội tụ đều thì $S'(x) = \sum u_n'(x)$ (đạo hàm từng số hạng).
  \end{enumerate}
\end{tinhchat}

\subsection{Chuỗi lũy thừa}

\begin{dinhnghia}
  \textbf{Chuỗi lũy thừa:}
  \[
    \sum_{n=0}^{\infty} a_n (x - x_0)^n = a_0 + a_1(x-x_0) + a_2(x-x_0)^2 + \cdots
  \]
\end{dinhnghia}

\begin{dinhly}[title={Định lý Abel}]
  Nếu chuỗi $\sum a_n x^n$ hội tụ tại $x = x_1 \ne 0$ thì nó hội tụ tuyệt đối với mọi $|x| < |x_1|$. Nếu phân kỳ tại $x_2$ thì phân kỳ với mọi $|x| > |x_2|$.
\end{dinhly}

\begin{dinhnghia}[title={Bán kính hội tụ}]
  Tồn tại $R \ge 0$ (hoặc $R = +\infty$) sao cho $\sum a_n x^n$:
  \begin{itemize}
    \item Hội tụ tuyệt đối khi $|x| < R$.
    \item Phân kỳ khi $|x| > R$.
    \item Khi $|x| = R$: cần xét riêng.
  \end{itemize}
  Khoảng $(-R, R)$ gọi là \textbf{khoảng hội tụ}.
\end{dinhnghia}

\begin{congthuc}[title={Công thức tính bán kính hội tụ}]
  \[
    R = \lim_{n \to \infty} \left|\frac{a_n}{a_{n+1}}\right| \quad\text{hoặc}\quad R = \frac{1}{\displaystyle\limsup_{n \to \infty} \sqrt[n]{|a_n|}} \quad\text{(Cauchy--Hadamard)}.
  \]
\end{congthuc}

\begin{tinhchat}[title={Tính chất chuỗi lũy thừa}]
  Trên khoảng hội tụ $(-R, R)$:
  \begin{enumerate}
    \item Tổng $S(x) = \sum a_n x^n$ liên tục.
    \item Có thể tích phân từng số hạng: $\displaystyle\int_0^x S(t)\,dt = \sum_{n=0}^{\infty} \frac{a_n}{n+1} x^{n+1}$.
    \item Có thể đạo hàm từng số hạng: $S'(x) = \sum_{n=1}^{\infty} n\,a_n x^{n-1}$.
    \item Chuỗi đạo hàm và chuỗi tích phân có cùng bán kính hội tụ $R$.
  \end{enumerate}
\end{tinhchat}

\begin{congthuc}[title={Khai triển Maclaurin thường gặp (bán kính hội tụ)}]
  \begin{align*}
    e^x &= \sum_{n=0}^{\infty} \frac{x^n}{n!} & (R = +\infty), \\
    \sin x &= \sum_{n=0}^{\infty} \frac{(-1)^n x^{2n+1}}{(2n+1)!} & (R = +\infty), \\
    \cos x &= \sum_{n=0}^{\infty} \frac{(-1)^n x^{2n}}{(2n)!} & (R = +\infty), \\
    \ln(1+x) &= \sum_{n=1}^{\infty} \frac{(-1)^{n-1} x^n}{n} & (R = 1), \\
    \frac{1}{1-x} &= \sum_{n=0}^{\infty} x^n & (R = 1), \\
    (1+x)^\alpha &= \sum_{n=0}^{\infty} \binom{\alpha}{n} x^n & (R = 1), \\
    \arctan x &= \sum_{n=0}^{\infty} \frac{(-1)^n x^{2n+1}}{2n+1} & (R = 1).
  \end{align*}
\end{congthuc}

\begin{phuongphap}[title={Ứng dụng tính gần đúng}]
  \begin{itemize}
    \item \textbf{Tính xấp xỉ hàm số:} Dùng tổng riêng $S_n(x)$ của chuỗi lũy thừa, sai số $|R_n| \le a_{n+1}$ (chuỗi đan dấu).
    \item \textbf{Tính tích phân:} Khi $\displaystyle\int f(x)\,dx$ không tính được bằng hàm sơ cấp, khai triển $f(x)$ thành chuỗi lũy thừa rồi tích phân từng số hạng.
  \end{itemize}
\end{phuongphap}

\begin{vidu}
  $\displaystyle\int_0^1 e^{-x^2}\,dx = \int_0^1 \sum_{n=0}^{\infty} \frac{(-1)^n x^{2n}}{n!}\,dx = \sum_{n=0}^{\infty} \frac{(-1)^n}{n!(2n+1)} = 1 - \frac{1}{3} + \frac{1}{10} - \frac{1}{42} + \cdots$
\end{vidu}

\subsection{Chuỗi Fourier}

\subsubsection{Chuỗi lượng giác}

\begin{dinhnghia}
  \textbf{Chuỗi lượng giác:}
  \[
    \frac{a_0}{2} + \sum_{n=1}^{\infty} \bigl(a_n \cos nx + b_n \sin nx\bigr).
  \]
  Hệ hàm $\{1, \cos x, \sin x, \cos 2x, \sin 2x, \ldots\}$ là hệ \textbf{trực giao} trên $[-\pi, \pi]$:
  \[
    \int_{-\pi}^{\pi} \cos mx \cdot \cos nx\,dx = \begin{cases} 0 & m \ne n, \\ \pi & m = n \ne 0, \end{cases}
  \]
  và tương tự cho $\sin$ và tích chéo $\sin \cdot \cos$.
\end{dinhnghia}

\subsubsection{Chuỗi Fourier}

\begin{dinhnghia}
  Cho $f(x)$ tuần hoàn chu kỳ $2\pi$, khả tích trên $[-\pi, \pi]$. \textbf{Chuỗi Fourier} của $f$:
  \[
    f(x) \sim \frac{a_0}{2} + \sum_{n=1}^{\infty} \bigl(a_n \cos nx + b_n \sin nx\bigr),
  \]
  với các \textbf{hệ số Fourier}:
  \[
    a_n = \frac{1}{\pi}\int_{-\pi}^{\pi} f(x)\cos nx\,dx, \qquad b_n = \frac{1}{\pi}\int_{-\pi}^{\pi} f(x)\sin nx\,dx.
  \]
\end{dinhnghia}

\begin{dinhly}[title={Điều kiện Dirichlet (điều kiện đủ hội tụ)}]
  Nếu $f$ tuần hoàn $2\pi$ và trên $[-\pi, \pi]$:
  \begin{enumerate}
    \item $f$ liên tục (trừ có thể hữu hạn điểm gián đoạn loại 1),
    \item $f$ đơn điệu từng khúc (hữu hạn khoảng đơn điệu),
  \end{enumerate}
  thì chuỗi Fourier hội tụ và:
  \[
    \frac{a_0}{2} + \sum_{n=1}^{\infty} (a_n \cos nx + b_n \sin nx) = \begin{cases} f(x) & \text{tại điểm liên tục,} \\[4pt] \dfrac{f(x^-) + f(x^+)}{2} & \text{tại điểm gián đoạn.} \end{cases}
  \]
\end{dinhly}

\subsubsection{Khai triển trên đoạn bất kỳ}

\begin{congthuc}[title={Chuỗi Fourier trên $[-l,\, l]$}]
  Với $f$ tuần hoàn chu kỳ $2l$:
  \[
    f(x) \sim \frac{a_0}{2} + \sum_{n=1}^{\infty} \left(a_n \cos\frac{n\pi x}{l} + b_n \sin\frac{n\pi x}{l}\right),
  \]
  \[
    a_n = \frac{1}{l}\int_{-l}^{l} f(x)\cos\frac{n\pi x}{l}\,dx, \qquad b_n = \frac{1}{l}\int_{-l}^{l} f(x)\sin\frac{n\pi x}{l}\,dx.
  \]
\end{congthuc}

\begin{phuongphap}[title={Khai triển hàm trên $[0, l]$}]
  \begin{itemize}
    \item \textbf{Khai triển chẵn} (chỉ chứa $\cos$): mở rộng $f$ thành hàm chẵn trên $[-l, l]$, khi đó $b_n = 0$.
    \item \textbf{Khai triển lẻ} (chỉ chứa $\sin$): mở rộng $f$ thành hàm lẻ trên $[-l, l]$, khi đó $a_n = 0$.
  \end{itemize}
\end{phuongphap}

\begin{dinhly}[title={Đẳng thức Parseval}]
  Nếu $f$ khả tích bình phương trên $[-\pi, \pi]$ thì:
  \[
    \frac{1}{\pi}\int_{-\pi}^{\pi} f^2(x)\,dx = \frac{a_0^2}{2} + \sum_{n=1}^{\infty} \bigl(a_n^2 + b_n^2\bigr).
  \]
  Đẳng thức này thể hiện sự bảo toàn năng lượng (bằng nhau giữa miền thời gian và miền tần số).
\end{dinhly}

%% ================================================================
%%  CHƯƠNG 2: PHƯƠNG TRÌNH VI PHÂN
%% ================================================================
\section{Phương trình vi phân}

\subsection{Phương trình vi phân cấp 1}

\subsubsection{Đại cương}

\begin{dinhnghia}
  \textbf{PTVP cấp 1} có dạng tổng quát $F(x, y, y') = 0$, hay dạng chính tắc $y' = f(x, y)$.

  \begin{itemize}
    \item \textbf{Nghiệm}: hàm $y = \varphi(x)$ thỏa mãn PT trên khoảng $(a, b)$.
    \item \textbf{Nghiệm tổng quát}: $\Phi(x, y, C) = 0$ chứa 1 hằng số tùy ý $C$.
    \item \textbf{Nghiệm riêng}: ứng với giá trị cụ thể của $C$.
    \item \textbf{Bài toán Cauchy}: tìm nghiệm thỏa điều kiện ban đầu $y(x_0) = y_0$.
  \end{itemize}
\end{dinhnghia}

\begin{dinhly}[title={Định lý Picard (tồn tại và duy nhất)}]
  Nếu $f(x,y)$ liên tục và thỏa điều kiện Lipschitz theo $y$ trong lân cận $(x_0, y_0)$ thì bài toán Cauchy $y' = f(x,y)$, $y(x_0) = y_0$ có nghiệm duy nhất.
\end{dinhly}

\subsubsection{Phương trình khuyết}

\begin{congthuc}
  \begin{itemize}
    \item \textbf{Khuyết $y$}: $F(x, y') = 0$ --- đặt $y' = p$, giải PT theo $x$ và $p$.
    \item \textbf{Khuyết $x$}: $F(y, y') = 0$ --- đặt $y' = p$, coi $p = p(y)$.
  \end{itemize}
\end{congthuc}

\subsubsection{PT biến số phân li}

\begin{dinhnghia}
  Dạng $M(x)\,dx + N(y)\,dy = 0$, hay $g(y)\,dy = f(x)\,dx$.
\end{dinhnghia}

\begin{phuongphap}
  Tích phân hai vế: $\displaystyle\int g(y)\,dy = \int f(x)\,dx + C$.
\end{phuongphap}

\subsubsection{PT thuần nhất}

\begin{dinhnghia}
  $y' = f\!\left(\dfrac{y}{x}\right)$, hay $M(x,y)\,dx + N(x,y)\,dy = 0$ với $M$, $N$ thuần nhất cùng bậc.
\end{dinhnghia}

\begin{phuongphap}
  Đặt $u = \dfrac{y}{x}$ ($y = ux$, $y' = u + xu'$), PT trở thành biến số phân li theo $x$ và $u$.
\end{phuongphap}

\subsubsection{PT tuyến tính cấp 1}

\begin{dinhnghia}
  $y' + P(x)\,y = Q(x)$.
  \begin{itemize}
    \item $Q(x) = 0$: PT tuyến tính \textbf{thuần nhất}.
    \item $Q(x) \ne 0$: PT tuyến tính \textbf{không thuần nhất}.
  \end{itemize}
\end{dinhnghia}

\begin{congthuc}[title={Nghiệm tổng quát}]
  \[
    y = e^{-\int P\,dx}\left(\int Q\,e^{\int P\,dx}\,dx + C\right).
  \]
\end{congthuc}

\begin{phuongphap}[title={Phương pháp biến thiên hằng số (Lagrange)}]
  \begin{enumerate}
    \item Giải PT thuần nhất $y' + Py = 0$: $y_{tn} = Ce^{-\int P\,dx}$.
    \item Thay $C = C(x)$, đặt $y = C(x)\,e^{-\int P\,dx}$ vào PT đầy đủ, tìm $C(x)$.
  \end{enumerate}
\end{phuongphap}

\subsubsection{PT Bernoulli}

\begin{dinhnghia}
  $y' + P(x)\,y = Q(x)\,y^\alpha$ \quad ($\alpha \ne 0, 1$).
\end{dinhnghia}

\begin{phuongphap}
  Chia hai vế cho $y^\alpha$, đặt $z = y^{1-\alpha}$, PT trở thành tuyến tính cấp 1 theo $z$:
  \[
    z' + (1-\alpha)P(x)\,z = (1-\alpha)Q(x).
  \]
\end{phuongphap}

\subsubsection{PT vi phân toàn phần}

\begin{dinhnghia}
  $P(x,y)\,dx + Q(x,y)\,dy = 0$ là \textbf{toàn phần} nếu $\dfrac{\partial P}{\partial y} = \dfrac{\partial Q}{\partial x}$.

  Khi đó tồn tại $U(x,y)$ sao cho $dU = P\,dx + Q\,dy$, và nghiệm: $U(x,y) = C$.
\end{dinhnghia}

\begin{phuongphap}[title={Tìm $U$}]
  $U(x,y) = \displaystyle\int_{x_0}^{x} P(t, y_0)\,dt + \int_{y_0}^{y} Q(x, s)\,ds$, hoặc:

  $U = \displaystyle\int P\,dx + \varphi(y)$, rồi xác định $\varphi(y)$ từ $U'_y = Q$.
\end{phuongphap}

\begin{phuongphap}[title={Thừa số tích phân}]
  Nếu PT chưa toàn phần, nhân với $\mu(x,y)$ sao cho $\mu P\,dx + \mu Q\,dy = 0$ trở thành toàn phần:
  \begin{itemize}
    \item Nếu $\dfrac{1}{Q}\!\left(\dfrac{\partial P}{\partial y} - \dfrac{\partial Q}{\partial x}\right)$ chỉ phụ thuộc $x$: $\mu = e^{\int (\ldots)\,dx}$.
    \item Nếu $\dfrac{1}{P}\!\left(\dfrac{\partial Q}{\partial x} - \dfrac{\partial P}{\partial y}\right)$ chỉ phụ thuộc $y$: $\mu = e^{\int (\ldots)\,dy}$.
  \end{itemize}
\end{phuongphap}

\subsubsection{PT Clairaut}

\begin{dinhnghia}
  $y = xy' + \psi(y')$.
\end{dinhnghia}

\begin{congthuc}
  \begin{itemize}
    \item \textbf{Nghiệm tổng quát} (họ đường thẳng): $y = Cx + \psi(C)$.
    \item \textbf{Nghiệm kỳ dị} (hình bao): khử $p$ từ $\begin{cases} y = xp + \psi(p), \\ 0 = x + \psi'(p). \end{cases}$
  \end{itemize}
\end{congthuc}

\subsubsection{PT Lagrange}

\begin{dinhnghia}
  $y = x\,\varphi(y') + \psi(y')$ \quad ($\varphi(p) \ne p$ để không phải Clairaut).
\end{dinhnghia}

\begin{phuongphap}
  Đặt $y' = p$, đạo hàm hai vế theo $x$, thu PT tuyến tính cấp 1 theo $x = x(p)$.
\end{phuongphap}

\subsubsection{Quỹ đạo trực giao}

\begin{dinhnghia}
  Cho họ đường $\Phi(x, y, C) = 0$. \textbf{Quỹ đạo trực giao} là đường cắt vuông góc mọi đường trong họ.
\end{dinhnghia}

\begin{phuongphap}
  \begin{enumerate}
    \item Tìm PTVP của họ: khử $C$ từ $\Phi = 0$ và $\Phi'_x + \Phi'_y \cdot y' = 0$, được $y' = f(x,y)$.
    \item PTVP quỹ đạo trực giao: $y' = -\dfrac{1}{f(x,y)}$.
    \item Giải PTVP mới.
  \end{enumerate}
\end{phuongphap}

\subsection{Phương trình vi phân cấp 2}

\subsubsection{Đại cương}

\begin{dinhnghia}
  Dạng tổng quát: $F(x, y, y', y'') = 0$, hay $y'' = f(x, y, y')$.

  Nghiệm tổng quát chứa 2 hằng số tùy ý. Bài toán Cauchy: $y(x_0) = y_0$, $y'(x_0) = y_0'$.
\end{dinhnghia}

\subsubsection{PT cấp 2 khuyết}

\begin{phuongphap}
  \begin{itemize}
    \item \textbf{Khuyết $y$ và $y'$}: $y'' = f(x)$ --- tích phân hai lần.
    \item \textbf{Khuyết $y$}: $y'' = f(x, y')$ --- đặt $p = y'$, được PT cấp 1: $p' = f(x, p)$.
    \item \textbf{Khuyết $x$}: $y'' = f(y, y')$ --- đặt $p = y'$, $y'' = p\,\dfrac{dp}{dy}$, được $p\,\dfrac{dp}{dy} = f(y, p)$.
  \end{itemize}
\end{phuongphap}

\subsubsection{PT tuyến tính cấp 2}

\begin{dinhnghia}
  $y'' + p(x)\,y' + q(x)\,y = f(x)$.
  \begin{itemize}
    \item $f(x) = 0$: PT tuyến tính \textbf{thuần nhất} (PTTTTN).
    \item $f(x) \ne 0$: PT tuyến tính \textbf{không thuần nhất}.
  \end{itemize}
\end{dinhnghia}

\begin{dinhly}[title={Cấu trúc nghiệm}]
  \begin{itemize}
    \item Nghiệm tổng quát PTTTTN: $y_{tn} = C_1 y_1 + C_2 y_2$, với $y_1, y_2$ là hệ nghiệm cơ bản (độc lập tuyến tính, tức Wronskian $W(y_1, y_2) \ne 0$).
    \item Nghiệm tổng quát PT không thuần nhất: $y = y_{tn} + y_{rn}$, trong đó $y_{rn}$ là một nghiệm riêng.
  \end{itemize}
\end{dinhly}

\begin{dinhnghia}[title={Wronskian}]
  \[
    W(y_1, y_2) = \begin{vmatrix} y_1 & y_2 \\ y_1' & y_2' \end{vmatrix} = y_1 y_2' - y_2 y_1'.
  \]
  $y_1, y_2$ độc lập tuyến tính $\Leftrightarrow$ $W \ne 0$.
\end{dinhnghia}

\begin{phuongphap}[title={Biến thiên hằng số (tìm nghiệm riêng)}]
  Đặt $y_{rn} = C_1(x)\,y_1 + C_2(x)\,y_2$, giải hệ:
  \[
    \begin{cases} C_1' y_1 + C_2' y_2 = 0, \\ C_1' y_1' + C_2' y_2' = f(x). \end{cases}
  \]
\end{phuongphap}

\subsubsection{PT tuyến tính cấp 2 hệ số không đổi}

\begin{dinhnghia}
  $y'' + py' + qy = f(x)$ \quad ($p, q$ là hằng số).
\end{dinhnghia}

\begin{congthuc}[title={Giải PT thuần nhất}]
  PT đặc trưng: $k^2 + pk + q = 0$.
  \begin{itemize}
    \item Hai nghiệm thực phân biệt $k_1 \ne k_2$: $y_{tn} = C_1 e^{k_1 x} + C_2 e^{k_2 x}$.
    \item Nghiệm kép $k_1 = k_2 = k$: $y_{tn} = (C_1 + C_2 x)e^{kx}$.
    \item Nghiệm phức $k = \alpha \pm \beta i$: $y_{tn} = e^{\alpha x}(C_1\cos\beta x + C_2\sin\beta x)$.
  \end{itemize}
\end{congthuc}

\begin{phuongphap}[title={Nghiệm riêng --- Phương pháp hệ số bất định}]
  \begin{itemize}
    \item $f(x) = e^{\alpha x} P_n(x)$: tìm $y_{rn} = x^s e^{\alpha x} Q_n(x)$, $s$ = bội của $\alpha$ trong PT đặc trưng.
    \item $f(x) = e^{\alpha x}\bigl[P_n(x)\cos\beta x + R_m(x)\sin\beta x\bigr]$:
      tìm $y_{rn} = x^s e^{\alpha x}\bigl[\tilde{P}_l(x)\cos\beta x + \tilde{R}_l(x)\sin\beta x\bigr]$, $l = \max(n,m)$, $s$ = bội của $\alpha + \beta i$.
  \end{itemize}
\end{phuongphap}

\subsubsection{PT Euler}

\begin{dinhnghia}
  $x^2 y'' + axy' + by = f(x)$ \quad ($a, b$ hằng số).
\end{dinhnghia}

\begin{phuongphap}
  Đặt $x = e^t$ (khi $x > 0$): $xy' = \dot{y}$, $x^2 y'' = \ddot{y} - \dot{y}$, PT trở thành PT hệ số không đổi theo $t$.
\end{phuongphap}

\subsubsection{PT dao động}

\begin{congthuc}
  \begin{itemize}
    \item \textbf{Dao động tự do}: $y'' + \omega^2 y = 0$ $\Rightarrow$ $y = A\cos(\omega x + \varphi)$.
    \item \textbf{Dao động tắt dần}: $y'' + 2\beta y' + \omega^2 y = 0$ ($\beta < \omega$) $\Rightarrow$ $y = Ae^{-\beta x}\cos(\omega_1 x + \varphi)$, $\omega_1 = \sqrt{\omega^2 - \beta^2}$.
    \item \textbf{Dao động cưỡng bức}: $y'' + 2\beta y' + \omega^2 y = F_0\cos\Omega x$.
    \item \textbf{Cộng hưởng}: xảy ra khi $\Omega = \omega$ (dao động tự do) hoặc $\Omega = \sqrt{\omega^2 - 2\beta^2}$ (có cản).
  \end{itemize}
\end{congthuc}

\subsubsection{Nghiệm khai triển thành chuỗi lũy thừa}

\begin{phuongphap}
  Tìm nghiệm dạng $y = \displaystyle\sum_{n=0}^{\infty} a_n x^n$: thế vào PT, so sánh hệ số cùng bậc, tìm công thức truy hồi cho $a_n$.
\end{phuongphap}

\subsection{Hệ phương trình vi phân}

\subsubsection{Đại cương}

\begin{dinhnghia}
  Hệ PTVP cấp 1 dạng chuẩn tắc:
  \[
    \begin{cases}
      x_1' = f_1(t, x_1, x_2, \ldots, x_n), \\
      x_2' = f_2(t, x_1, x_2, \ldots, x_n), \\
      \vdots \\
      x_n' = f_n(t, x_1, x_2, \ldots, x_n).
    \end{cases}
  \]
  Dạng ma trận: $\mathbf{x}' = \mathbf{f}(t, \mathbf{x})$.
\end{dinhnghia}

\begin{luuy}
  Mọi PTVP cấp $n$ đều chuyển được về hệ cấp 1: đặt $x_1 = y$, $x_2 = y'$, \ldots, $x_n = y^{(n-1)}$.
\end{luuy}

\subsubsection{Cách giải hệ PTVP}

\begin{phuongphap}[title={Phương pháp khử}]
  Từ hệ, khử dần các ẩn để được 1 PT cấp cao theo 1 ẩn, giải rồi thế ngược lại.
\end{phuongphap}

\subsubsection{Hệ tuyến tính thuần nhất hệ số không đổi}

\begin{dinhnghia}
  $\mathbf{x}' = A\,\mathbf{x}$, với $A$ là ma trận hằng $n \times n$.
\end{dinhnghia}

\begin{phuongphap}[title={Phương pháp trị riêng}]
  \begin{enumerate}
    \item Tìm trị riêng $\lambda$ từ PT đặc trưng $\det(A - \lambda I) = 0$.
    \item Với mỗi $\lambda_i$, tìm vectơ riêng $\mathbf{v}_i$: $(A - \lambda_i I)\mathbf{v}_i = \mathbf{0}$.
    \item Nghiệm cơ bản: $\mathbf{x}_i(t) = e^{\lambda_i t}\,\mathbf{v}_i$.
    \item Nghiệm tổng quát: $\mathbf{x}(t) = C_1\,\mathbf{x}_1(t) + C_2\,\mathbf{x}_2(t) + \cdots + C_n\,\mathbf{x}_n(t)$.
  \end{enumerate}
\end{phuongphap}

\begin{congthuc}[title={Các trường hợp (hệ $2 \times 2$)}]
  PT đặc trưng $\lambda^2 - \text{tr}(A)\,\lambda + \det(A) = 0$:
  \begin{itemize}
    \item \textbf{Hai trị riêng thực phân biệt} $\lambda_1 \ne \lambda_2$:
      $\mathbf{x} = C_1 e^{\lambda_1 t}\mathbf{v}_1 + C_2 e^{\lambda_2 t}\mathbf{v}_2$.
    \item \textbf{Trị riêng kép} $\lambda$, thiếu vectơ riêng:
      $\mathbf{x} = C_1 e^{\lambda t}\mathbf{v} + C_2 e^{\lambda t}(t\mathbf{v} + \mathbf{w})$, trong đó $(A - \lambda I)\mathbf{w} = \mathbf{v}$.
    \item \textbf{Trị riêng phức} $\lambda = \alpha \pm \beta i$, vectơ riêng $\mathbf{v} = \mathbf{p} + i\mathbf{q}$:
      \begin{align*}
        \mathbf{x} &= C_1 e^{\alpha t}(\mathbf{p}\cos\beta t - \mathbf{q}\sin\beta t) + C_2 e^{\alpha t}(\mathbf{p}\sin\beta t + \mathbf{q}\cos\beta t).
      \end{align*}
  \end{itemize}
\end{congthuc}

\begin{phuongphap}[title={Ma trận mũ}]
  Nghiệm tổng quát: $\mathbf{x}(t) = e^{At}\,\mathbf{x}(0)$, với $e^{At} = I + At + \dfrac{(At)^2}{2!} + \cdots$

  Hệ không thuần nhất $\mathbf{x}' = A\mathbf{x} + \mathbf{b}(t)$:
  \[
    \mathbf{x}(t) = e^{At}\mathbf{x}(0) + \int_0^t e^{A(t-s)}\mathbf{b}(s)\,ds.
  \]
\end{phuongphap}

\end{document}